% !TEX root = ../../tree_notes.tex
% =====================================================================
%  Chapter 22 -- Tree Isomorphism and Hashing
%  Status: OUTLINE (not written yet). See structure/chapter.txt.
% =====================================================================
\bookchapter{Tree Isomorphism and Hashing}{c22}
\chapterinfo{CF 1800--2300}{Outline}

\plannedtopic{Rooted Tree Isomorphism and AHU}{9}{%
  \item Two rooted trees are isomorphic if a bijection preserves the parent--child relation
  \item The AHU algorithm: canonical labels bottom-up by sorted child labels
  \item Assigning ids to sorted child-label tuples (map of vectors)
  \item $O(n\log n)$ with sorting, $O(n)$ with radix-style bucketing
  \item Handling unrooted trees: try the centre(s) from \chref{c09}
}

\plannedtopic{Canonical Forms and the Tree Centre}{8}{%
  \item The centre of a tree (one node or an edge) and why it is the right root
  \item Canonical form / encoding of a tree as a balanced-parentheses string
  \item Comparing two trees by their canonical strings
  \item Rooting at the centre makes isomorphism canonical
}

\plannedtopic{Tree Hashing}{9}{%
  \item A randomized polynomial hash of a subtree: combine child hashes commutatively
  \item XOR/add of a random value per child, or a sorted-tuple hash
  \item Collision risk and using 64-bit / double moduli
  \item Comparing subtrees of one big tree in $O(1)$ after $O(n)$ preprocessing
  \item Hashing vs AHU: when a hash is enough
}

\plannedtopic{Counting Distinct Subtrees and Shapes}{8}{%
  \item Number of distinct rooted subtree shapes in a tree
  \item Counting unlabeled trees and the link to Catalan/Cayley (\chref{c01}, \chref{c17})
  \item Distinct-substring-style counting on the tree
  \item Isomorphism classes of forests
}

\plannedtopic{Isomorphism Practice}{7}{%
  \item Subtree of another tree via hashing
  \item Find all isomorphic subtrees / repeated structure
  \item Counting non-isomorphic trees with $n$ nodes (small $n$ tables)
}

\closechapter
